\section{Related Work}\label{sec:related-work}

\ly{TODO: There are tons of related work not mentioned here yet. On lazy
  computation cost:
  \cite{iris-credit,iris-thunk,liquidate,liquid-okasaki,danielsson-08,foner2018,first-order-lazy,lazy-amortized-testing}.
  There are also many papers on non-lazy computation costs.}

\paragraph{Lazy semantics.}
There are many semantics for lazy programming languages. The best known of these
is the natural semantics of Launchbury~\cite{launchbury1993}. Clairvoyance
semantics was first studied by Hackett and Hutton~\cite{clairvoyant}, who proved
its equivalence with Launchbury's semantics. Li \emph{et al.} developed a model
of this semantics in Rocq, along with a framework for automating proofs about
the behavior of programs in it~\cite{forking-paths}. Demand semantics was first
studied by Bjerner and Holmstr{\"{o}}m~\cite{bjerner1989}. Their semantics was
untyped; a typed version, essentially the same as the demand semantics in the
present work, was first published by Xia \emph{et
  al.}~\cite{demand-semantics}. Xia \emph{et al.} were also the first to
demonstrate that demand semantics is equivalent to other semantics for lazy
languages; in particular, they proved the cost existence and cost minimality
theorems, which, together, amount to equivalence with clairvoyance semantics.

\paragraph{Lazy cost analysis.}
This work grew out of an effort to formally verify the performance of lazy
programs: the original motivation was to find a ``source of truth'' against
which to judge the correctness of hand-written demand
terms. Okasaki's~\cite{purely-functional} analysis techniques are foundational
to this area of research. Formal developments have been carried out in Rocq, particularly via the Iris framework~\cite{iris-credit,iris-thunk}; in Liquid Haskell~\cite{liquidate,liquid-okasaki}; and in Agda~\cite{danielsson-08}. There is also work on \emph{testing} the performance on lazy data structure; \eg\ by Lorenzen~\cite{lazy-amortized-testing}.

\paragraph{Lazy performance.}
Beyond just analyzing costs, there is great value in improving the performance
of lazy programs. The work in this area is far too vast to summarize succinctly;
see Jones and Salkild~\cite{DBLP:conf/fpca/JonesS89}, Jones and
Partain~\cite{strictness-analysis}, Massey and
Tick~\cite{DBLP:conf/ifipPACT/MasseyT94}, and Lorenzen \textit{et
  al.}~\cite{10.1145/3747530}.

\paragraph{Cost analysis.}

\paragraph{Functional logic programming.}
We implemented our examples in Curry; however, most of the techniques in this
paper are not specific to Curry, and could have been implemented in other
functional logic languages such as Verse~\cite{verse} or
mercury~\cite{mercury}. Because \(\LL\) uses eager evaluation, we could have
even translated to a pure logic language; \eg\ Prolog~\cite{EagerFLPIsLogic}.
~\cite{DBLP:journals/jflp/BrasselHH04}. However, the search tree traversal employed in this work is Curry-specific~\cite{DBLP:journals/jflp/BrasselHH04}.
